\section{Models and Assumptions}
\begin{frame}{\hspace{3cm}}
\centering

\vspace{1.5cm}

\begin{beamercolorbox}[sep=20pt,center,rounded=true,shadow=true]{title}
    {\fontsize{20}{24}\selectfont \textbf{Models and assumption}}
\end{beamercolorbox}

\vspace{0.8cm}


\vfill

\end{frame}
%------------------------------------------------

\begin{frame}{Model Framework}
Suppose that we observe $n$ i.i.d.\ $d$-dimensional observations 
$X_i = (X_{1i}, \ldots, X_{di})$, \quad $i = 1, \ldots, n$.
\begin{itemize}
    \item Under an ideal physical model ($H_0$):
    \begin{itemize}
        \item Observations $X_i$ have density:
        \[
        f(x_i, \theta)
        \]
        \item $\theta = (\theta_1, \dots, \theta_k) \in \mathbb{R}^k$
        \item $\theta$ is unknown and estimated from data
    \end{itemize}

    \vspace{0.3cm}

    \item Under less ideal conditions ($H_A$):
    \begin{itemize}
        \item True density is:
        \[
        f(x_i)
        \]
        \item No parametric form assumed
        \item Must be estimated from the data
    \end{itemize}

    \vspace{0.3cm}

    \item Due to generality of the procedure, several \textbf{regularity and rate assumptions} are required
\end{itemize}
\end{frame}

%------------------------------------------------

\begin{frame}{Likelihood and Estimation}
\begin{itemize}
    \item Parameter $\theta_0$ estimated using MLE:
    \[
    \prod_{i=1}^n f(x_i,\hat\theta) = \max_{\theta} \prod_{i=1}^n f(x_i,\theta)
    \]

    \item Nonparametric estimator (kernel density):
    \[
    \hat{f}(z) = \frac{1}{n b^d} \sum K\left(\frac{x_i - z}{b}\right)
   \]

    \item Semiparametric estimator:
    \[
    g(x,\pi) = \pi f(x,\hat\theta) + (1-\pi)\hat{f}(x), \quad 0 \leq \pi \leq 1
    \]
\end{itemize}
\end{frame}

%------------------------------------------------

\begin{frame}{Assumptions on MLE}
\begin{itemize}
    \item \textbf{Consistency and rate of convergence:}

    \[
    n \mathbb{E}\|\hat{\theta} - \theta_0\|^2 = O(1) \quad \text{under } H_0
    \]

    \[
    n \mathbb{E}\|\hat{\theta} - \theta_A\|^2 = O(1) \quad \text{under } H_A
    \]

    \item Interpretation:
    \begin{itemize}
        \item $\hat{\theta}$ converges, at $\sqrt{n}$ rate, to $\theta_0$ under $H_0$ and to $\theta_A$ under $H_A$.
    \end{itemize}
\end{itemize}
\end{frame}

%------------------------------------------------

\begin{frame}{Likelihood Ratio Assumptions}
Under $H_0$,
\begin{equation}
\sqrt{n}\,\mathbb{E}\{\tfrac{\sqrt{f(X_i,\hat\theta)} - \sqrt{f(X_i,\theta_0)}}{\sqrt{f(X_i,\hat\theta)}} \}^2 = O(1).
\tag{a}
\end{equation}

Under $H_A$,
\begin{equation}
\sqrt{n}\,\mathbb{E}\{\tfrac{\sqrt{f(X_i,\hat\theta)} - \sqrt{f(X_i,\theta_A)}}{\sqrt{f(X_i,\hat\theta)}} \}^2 = O(1).
\tag{b}
\end{equation}

As $n \to \infty$,
\begin{equation}
f(X_i,\hat\theta) \xrightarrow{L_2} f(X_i,\theta_0) \quad \text{under } H_0, 
\qquad 
f(X_i,\hat\theta) \xrightarrow{L_2} f(X_i,\theta_A) \quad \text{under } H_A.
\tag{c}
\end{equation}

Assumptions (a) and (b) require the expected likelihood ratios to converge to $1$ at the usual rates; (c) requires $L_2$ convergence of $f(x,\theta)$ to its limit.

In addition, we require conditions on the densities: $f(x)$ and $f(x,\theta)$ have compact support and
\[
\inf_x f(x) > \varepsilon, 
\qquad 
\inf_{x,\theta} f(x,\theta) > \varepsilon,
\]
where $\varepsilon$ is known. In other words, we require that the underlying density has compact support and is bounded below.
\end{frame}

%------------------------------------------------

\begin{frame}{Regularity Conditions }
\begin{itemize}
    \item Smoothness:
    
    $f(x,\theta)$ is twice continuously differentiable in each argument. 
    
    The derivative can be taken under the integral sign.


    \item Boundedness:
    \[
    \sup_x |f(x)| < M
    \]

    \end{itemize}
\end{frame}

%------------------------------------------------

\begin{frame}{Kernel Density Assumptions}
To provide some protection against the possibility that the data come from $H_A$ we use a common nonparametric kernel density estimate
\begin{itemize}
    \item Kernel estimator:
    \[
    \hat{f}(z) = \frac{1}{n b^d} \sum K\left(\frac{x_i - z}{b}\right)
    \]
where z is the point at which the density is to be estimated, K(x) is a specified density on $R^d$ and is a continuous kernel with compact support
    \item Conditions on bandwidth:
    \[
    b \to 0, \quad n b^{2d} \to \infty
    \]

    \item Ensures:
    \begin{itemize}
        \item Consistency
        \item Almost sure convergence
    \end{itemize}

    \item . It has been shown (Devroye 1983) that $\hat f(x)$ is an $L_1$ consistent estimate of $f(x)$.
\end{itemize}
\end{frame}

%------------------------------------------------

\begin{frame}{Non-parametric Rate Assumptions }
\begin{itemize}
    \item Mean square convergence:
    \[
    \mathbb{E}\left[
    \frac{\hat{f}(X_i) - f(X_i)}{f(X_i)}
    \right]^2 = O(a_n^2)
    \]

    \item Where:
    \begin{itemize}
    \item $a_n$ related to bandwidth b
        \item $a_n \to 0$
        \item $a_n n^{\alpha/2} \to \infty$ for some $\alpha$ between 0 and 1
        \item $f(x_i)$ is $f(x_i, \theta_0)$ when $H_0$ is true and $f(x_i)$ is the true density when $H_A$ is true.
    \end{itemize}
\end{itemize}
\end{frame}

%------------------------------------------------

\begin{frame}{Cross-Product Conditions}
\begin{itemize}
    \item Under $H_0$:
    \[
    \int [\hat{f}(x_i)-f(x_i, \theta_0)][\hat{f}(x_j)-f(x_j,\theta_0)] dx_i dx_j
    = O_p\left(\frac{a_n^2}{n}\right)\quad \text{for}\quad i\neq j
    \]

    \item Under $H_A$:
    \[
    \int [\hat{f}(x_i)-f(x_i)][\hat{f}(x_j)-f(x_j)] dx_i dx_j
    = O_p\left(\frac{a_n^2}{n}\right)\quad \text{for}\quad i \neq j 
    \]

    \item Interpretation:
    \begin{itemize}
        \item Controls dependence between estimation errors
    \end{itemize}
\end{itemize}
\end{frame}

%------------------------------------------------

\begin{frame}{Key Idea of the Model}
\begin{itemize}
    \item Final estimator:convex combination of the parametric and non-parametric estimates
    \[
    g(x,\pi) = \pi f(x,\hat{\theta}) + (1-\pi)\hat{f}(x)
    \]

    \item The parameter $\pi$ is unknown and its estimate $\hat{\pi}$ is obtained from the data, so that $g(x,\hat\pi)$ is the actual density estimate. 

Let
\[
\tilde L(\mathbf{x},\pi) = \prod_{i=1}^n g(x_i,\pi)
\]
denote the pseudolikelihood function from which we choose a pseudomaximum likelihood estimate $\hat{\pi}$ for $\pi$ that satisfies
\[
\tilde L(\mathbf{X},\hat{\pi}) = \prod_{i=1}^n g(x_i,\hat{\pi})
= \max_{\pi} \prod_{i=1}^n g(x_i,\pi).
\]

    \item Interpretation:
    When $H_0$ holds the chosen density should be close to $f(x,\theta_0)$, and when $H_A$ holds the density should be close to $f(x)$.
\end{itemize}
\end{frame}
